\begin{lemma}\label{lem:nonic81-Hecke-finite-support}
Let $N\geq5$ be odd, and suppose that $J_1(N)(\mathbb Q)$ is finite
and $J_1(N)(\mathbb Q)[2]=0$. If
\[
 \operatorname{gon}_{\mathbb C}X_1(N)>3d,
\]
then $X_1(N)$ has no rational effective noncuspidal divisor of degree
at most $d$.
\end{lemma}
\begin{proof}
Put $X=X_1(N)$, $J=J_1(N)$, and write $\delta=\langle2\rangle$.
We use the correspondence convention
\[
 T_2(E,P)=\sum_{C\subset E,\,\#C=2}(E/C,\phi_C(P)).
\]
Reduction at $2$ is injective on $J(\mathbb Q)$.
Indeed, the kernel on torsion is killed by $2$: in the formal group
at the identity, multiplication by $2$ maps parameters in
$2\mathbb Z_2$ to $4\mathbb Z_2$, and the latter subgroup is
torsion-free by the formal logarithm. The assumption on rational
$2$-torsion therefore gives injectivity.

The Eichler--Shimura relation on the special fibre is
\[
 F^2-T_2F+2\delta=0.
\]
On $J(\mathbb F_2)$, Frobenius acts as the identity. Consequently
$A=T_2-1-2\delta$ kills the reduction of $J(\mathbb Q)$, and injectivity
gives $A J(\mathbb Q)=0$.

Let $c$ be the rational cusp with Tate level structure
$(\operatorname{Tate}(q),q^{1/N})$. We have the literal divisor identity
\[
 T_2c=c+2\delta c.                                      \tag{*}
\]
To see this, quotienting the Tate curve by $\mu_2$ sends
$(q,z)$ to $(q^2,z^2)$ and preserves the component label $1/N$.
The other two cyclic subgroups give a Tate parameter $q'$ with
$(q')^2=q$ and preserve $z$. Their component label is $2/N$,
so both specialize to $\delta c$. This also fixes the diamond
convention in the Eichler--Shimura relation above.

Suppose that $D$ is a rational effective noncuspidal divisor of
degree $e\leq d$. Applying $A$ to $[D-ec]$ and using $(*)$, we obtain
\[
 T_2D\sim D+2\delta D.
\]
Both divisors are effective of degree $3e$. If their difference were
nonzero, a function with this divisor would have degree at most
$3e<\operatorname{gon}_{\mathbb C}X$. Hence
\[
 T_2D=D+2\delta D                                      \tag{**}
\]
as divisors.

The finite set of $j$-invariants in the support of $D$ is therefore
closed under every cyclic $2$-isogeny. This is impossible.
For an elliptic curve $E_\tau=\mathbb C/(\mathbb Z+\tau\mathbb Z)$,
with $\operatorname{Im}(\tau)>0$, multiplication by $2$ induces a
cyclic $2$-isogeny $E_\tau\longrightarrow E_{2\tau}$.
Iterating gives $j(2^k\tau)$ in the same finite set. However,
$j(2^k\tau)$ is unbounded as $k\to\infty$, by the expansion
$j(q)=q^{-1}+744+O(q)$. This contradicts $(**)$.\qedhere
\end{proof}
